\section{Taller resuelto: Un límite de CG con margen mínimo}

\noindent\textbf{Nivel:} Intermedio.

\apartado{1. Enunciado}{Con aw=5.0, at=4.3, St/S=0.20, eta=0.93, $epsilon_{alpha}$=0.33, hac=0.25, ht=3.3 y $Cm_{alpha,f}$=0.09, determine el CG más retrasado que conserva margen estático mínimo 0.08.}

\apartado{2. Datos}{Posiciones adimensionales hacia atrás; pendientes por radián; mando fijo.}

\apartado{3. Incógnitas}{Las magnitudes y comprobaciones solicitadas en el caso.}

\apartado{4. Sistema de referencia}{Momento de cabeceo positivo morro arriba; referencia de posición común.}

\apartado{5. Hipótesis físicas}{Modelo lineal con punto neutro independiente de CG dentro de las hipótesis declaradas.}

\apartado{6. Modelo matemático}{Un límite de CG con margen mínimo}

\apartado{7. Ecuaciones simbólicas}{Se parte de los balances y definiciones siguientes.}

\begin{equation}K_t=\eta_t(S_t/S)a_t(1-\varepsilon_\alpha)\end{equation}

\begin{equation}h_n=(a_wh_{ac}+K_th_t-C_{m\alpha,f})/(a_w+K_t)\end{equation}

\begin{equation}h_{max}=h_n-SM_{min}\end{equation}

\apartado{8. Desarrollo algebraico}{Las operaciones se realizan en este orden, conservando las condiciones del modelo.}

\begin{enumerate}

\item Se calcula punto neutro conservando sustentación de cola en la pendiente total.

\item Se resta el margen exigido para fijar el límite retrasado.

\item Se evalúa la pendiente de momento en ese límite para comprobar el requisito.

\end{enumerate}

\apartado{9. Sustitución numérica}{Se usan las magnitudes del apartado de datos en SI; los ángulos se convierten a radianes antes de evaluar funciones trigonométricas. Las densidades indicadas como ISA se obtienen de la capa correspondiente.}

\noindent Punto neutro:\quad $0.5289791696547569\simeq 0.5289792\,$.\par

\noindent CG máximo:\quad $0.4489791696547569\simeq 0.4489792\,$.\par

\noindent $Cm_{alpha}$ en límite:\quad $-0.44286928\simeq -0.4428693\,\mathrm{rad^{-1}}$.\par

\apartado{10. Resultados con unidades}{Los resultados pertenecen al modelo y al intervalo de datos declarados.}

\begin{center}\begin{tabular}{lr}\toprule Magnitud&Resultado\\\midrule

Punto neutro & $0.5289792\,$\\

CG máximo & $0.4489792\,$\\

$Cm_{alpha}$ en límite & $-0.4428693\,\mathrm{rad^{-1}}$\\

\bottomrule\end{tabular}\end{center}

\apartado{11. Interpretación}{El requisito fija una frontera estática y no determina por sí solo límites de mando, estructura o certificación.}

\apartado{12. Comprobación independiente y alcance}{Sustituir el CG límite en la expresión directa de ala, cola y fuselaje y recuperar -a por 0.08.}

\apartado{13. Limitaciones}{El punto neutro puede cambiar con Mach, potencia, configuración o mando libre.}

\par\smallskip\noindent{\footnotesize\textit{Datos y solución originales. Las sustituciones registradas se reproducen con SymPy y Studium; eso no establece validación experimental ni revisión experta.}}
